\documentclass[11pt]{article}
\usepackage[margin=0.82in]{geometry}
\usepackage{amsmath,amssymb,amsthm,graphicx,microtype}
\usepackage[hidelinks]{hyperref}
\setlength{\parindent}{0pt}
\setlength{\parskip}{0.5em}
\newtheorem{theorem}{Theorem}
\newtheorem{remark}{Remark}
\newcommand{\OO}{\mathcal O}
\newcommand{\C}{\mathbb C}
\newcommand{\cl}{\overline}
\begin{document}

\begin{center}
{\Large\bfseries A Thomson--Brennan Theorem for Balanced Finite Covers}

\medskip
{\large A partial solution to the Open Problem of arXiv:2112.08504}

\smallskip
Lane 19, agent\_lane\_19, model GPT5.6
\end{center}

\textbf{Status.} \texttt{candidate\_partial\_solution\_likely\_valid}.

\textbf{Source problem.} Biswas--Putinar ask on PDF page 6 whether, for every
atomless positive measure compactly supported in a relatively compact
finite-connectivity domain $U$ of a finite-genus open Riemann surface,
$\OO(U)$ is dense in $L^2(\mu)$ exactly when there are no corresponding
analytic bounded point evaluations \cite{BP}.

\begin{center}
\includegraphics[width=0.94\textwidth]{figures/open_problem_crop.png}
\end{center}

\textbf{Definitions.} Let $\Omega\subset\C$ be a domain and let
$\pi:U\to\Omega$ be a finite unbranched holomorphic covering of degree $d$.
For an atomless finite positive Borel measure $\nu$ with compact support
$K\Subset\Omega$, define its \emph{balanced lift} $\mu$ by
\begin{equation}\label{balanced}
 \int_U F\,d\mu
 =\frac1d\int_\Omega\sum_{w\in\pi^{-1}(z)}F(w)\,d\nu(z)
 \qquad(F\ge0\text{ Borel}).
\end{equation}
Write $H^2(\Omega,\nu)=\cl{\OO(\Omega)}^{L^2(\nu)}$ and similarly on
$U$. A point is an analytic bounded point evaluation (ABPE) if evaluation is
bounded with a locally uniform bound throughout a neighborhood.

\textbf{Proof intuition.} A function on the cover is a holomorphic section of
the direct-image bundle $E=\pi_*\OO_U$, whose fiber consists of functions on
the $d$-point fiber of $\pi$. The bundle is holomorphically trivial because the
base is an open Riemann surface. The balanced lift equips it with the ordinary
counting metric. In any global holomorphic frame that metric is uniformly
equivalent, on $K$, to the Euclidean metric. Thus mean approximation upstairs
is exactly $d$ independent copies of mean approximation downstairs, up to an
equivalent norm. Point evaluations pass through the same frame, so Brennan's
planar dichotomy completes the argument.

\begin{theorem}[balanced finite-cover case]\label{thm:main}
Let $\Omega\subset\C$ have finitely many complementary components, each of
positive diameter. Let $\pi:U\to\Omega$ be a finite unbranched holomorphic
covering, let $\nu$ be atomless and compactly supported in $\Omega$, and let
$\mu$ be its balanced lift \eqref{balanced}. Then:
\begin{align*}
 H^2(U,\mu)=L^2(\mu)
 &\quad\Longleftrightarrow\quad
 H^2(\Omega,\nu)=L^2(\nu),\\
 H^2(U,\mu)\text{ has ABPEs}
 &\quad\Longleftrightarrow\quad
 H^2(\Omega,\nu)\text{ has ABPEs}.
\end{align*}
Consequently $\OO(U)$ is dense in $L^2(\mu)$ if and only if it has no ABPEs.
\end{theorem}

\begin{proof}
The direct image $E=\pi_*\OO_U$ is a rank-$d$ holomorphic vector bundle on
$\Omega$: over an evenly covered disk, a section is simply a $d$-tuple of
holomorphic functions. Every holomorphic vector bundle on a noncompact
Riemann surface is holomorphically trivial \cite[Section 30, in particular
Theorem 30.4]{Forster}. Choose a global holomorphic frame
$s_1,\ldots,s_d$. Hence every $g\in\OO(U)$ has a unique representation
\begin{equation}\label{frame}
 g=\sum_{j=1}^d (a_j\circ\pi)s_j,
 \qquad a=(a_1,\ldots,a_d)\in\OO(\Omega)^d.
\end{equation}

Give $E_z$ the normalized counting metric
\[
 \langle v,u\rangle_z=\frac1d
 \sum_{w\in\pi^{-1}(z)}v(w)\overline{u(w)}.
\]
Let $G(z)=(\langle s_j,s_k\rangle_z)_{j,k}$. It is a continuous positive
definite matrix. Compactness of $K=\operatorname{supp}\nu$ gives constants
$0<c\le C<\infty$ such that
\begin{equation}\label{gram}
 c\lVert\xi\rVert_2^2\le \xi^*G(z)\xi
 \le C\lVert\xi\rVert_2^2
 \qquad(z\in K,\ \xi\in\C^d).
\end{equation}
Equations \eqref{balanced} and \eqref{frame} yield
\begin{equation}\label{norm}
 \lVert g\rVert_{L^2(\mu)}^2
 =\int_K a(z)^*G(z)a(z)\,d\nu(z).
\end{equation}

The same fiberwise correspondence identifies all of $L^2(\mu)$ with the
measurable $L^2$ sections of $E$. In the chosen frame, \eqref{gram} says that
its norm is equivalent to the standard norm on $L^2(\nu)^d$. Under this
identification $\OO(U)$ is exactly $\OO(\Omega)^d$. Therefore it is dense if
and only if $\OO(\Omega)$ is dense in $L^2(\nu)$, proving the first
equivalence.

Suppose next that $z$ is an ABPE downstairs. Each coefficient in
\eqref{frame} satisfies
$|a_j(z)|\le M_z\lVert a_j\rVert_{L^2(\nu)}$, with $M_z$ locally bounded.
The lower bound in \eqref{gram} controls every coefficient norm by
$c^{-1/2}\lVert g\rVert_{L^2(\mu)}$. Since the frame sections are locally
bounded, for every $w\in\pi^{-1}(z)$,
\[
 |g(w)|\le \sum_j|s_j(w)|\,|a_j(z)|
 \le C_w\lVert g\rVert_{L^2(\mu)},
\]
with locally bounded constants. Thus all points over an ABPE neighborhood are
ABPEs upstairs.

Conversely, if $w$ is an ABPE upstairs and $z=\pi(w)$, apply its estimate to
$g=a\circ\pi$. Formula \eqref{balanced} gives
$\lVert a\circ\pi\rVert_{L^2(\mu)}=\lVert a\rVert_{L^2(\nu)}$ and
$g(w)=a(z)$. Hence $z$ is an ABPE downstairs. This proves the second
equivalence.

Finally, Runge approximation identifies the $L^2(\nu)$-closure of
$\OO(\Omega)$ with the rational closure used in Brennan's theorem. Under the
stated finite-connectivity and positive-diameter hypotheses, Brennan's
theorem says that this closure is proper exactly when it has ABPEs
\cite{Brennan,BP}. Combining that dichotomy with the two equivalences proves
the last assertion.
\end{proof}

\begin{remark}[genuinely positive genus]
The theorem is not confined to coordinate charts. A connected double cover
of a four-connected planar domain with nontrivial monodromy around each of
its four boundary components is an unbranched bordered surface with four
boundary components and genus one: multiplicativity of Euler characteristic
gives $2-2g-4=2(2-4)=-4$, hence $g=1$. Such a cover can be placed relatively
compactly in a slightly larger cover, fitting the topology of the source
problem.
\end{remark}

\textbf{Scope and remaining obstruction.} The proof uses the balanced fiber
law essentially. An arbitrary measure disintegrates into a possibly
degenerate matrix weight in a frame of $\pi_*\OO_U$; then \eqref{gram} can
fail in the only direction needed to recover coefficient evaluations. The
theorem also does not treat branching that accumulates at the support. The
full source problem remains open.

\textbf{Novelty and verification.} Searches through 2026-08-13 used the exact
problem wording and combinations of ``Thomson/Brennan theorem'', ``finite
covering'', ``balanced measure'', ``vector bundle'', and ``matrix-valued
measure''. They found no matching statement. The identification
$\OO(U)=\Gamma(\Omega,\pi_*\OO_U)$, bundle triviality, compact Gram bounds,
both BPE implications, and the Euler-characteristic example were checked
independently in the accompanying verification report. Novelty confidence is
moderate pending specialist review.

\textbf{Human-review recommendation.} Check first that the global frame of
$\pi_*\OO_U$ induces the claimed measurable $L^2$ identification, then the
uniform lower Gram bound on $K$. Those two points make the reduction to
Brennan exact.

\begin{thebibliography}{9}
\bibitem{BP}
S. Biswas and M. Putinar,
\emph{Approximation in the mean on rational curves},
arXiv:2112.08504 (2021).

\bibitem{Brennan}
J. E. Brennan,
\emph{The structure of certain spaces of analytic functions},
Comput. Methods Funct. Theory 8 (2008), 625--640.

\bibitem{Forster}
O. Forster,
\emph{Lectures on Riemann Surfaces}, Graduate Texts in Mathematics 81,
Springer, corrected printing, 1999.
\end{thebibliography}

\end{document}

